% =============================================================
%  APOSTILA DE MACHINE LEARNING — TEORIA AVANÇADA
%  Compilar com: pdflatex apostila_ml.tex (2x para refs cruzadas)
% =============================================================
\documentclass[12pt,a4paper,twoside]{book}

% ----- Codificação e língua -----
\usepackage[utf8]{inputenc}
\usepackage[T1]{fontenc}
\usepackage[brazil]{babel}

% ----- Tipografia -----
\usepackage{lmodern}
\usepackage{microtype}

% ----- Geometria -----
\usepackage[top=2.5cm,bottom=2.5cm,left=3cm,right=2.5cm,
            headheight=15pt]{geometry}

% ----- Matemática -----
\usepackage{amsmath,amssymb,amsthm,mathtools}
\usepackage{bm}          % bold math
\usepackage{bbm}         % blackboard bold

% ----- Gráficos e cores -----
\usepackage{graphicx}
\usepackage[dvipsnames,svgnames,table]{xcolor}

% ----- Tabelas -----
\usepackage{booktabs,multirow,array,longtable}

% ----- Algoritmos -----
\usepackage[ruled,vlined,linesnumbered]{algorithm2e}

% ----- Caixas e destaques -----
\usepackage[most]{tcolorbox}
\tcbuselibrary{theorems,skins,breakable}

% ----- Cabeçalho/rodapé -----
\usepackage{fancyhdr}
\pagestyle{fancy}
\fancyhf{}
\fancyhead[LE]{\small\leftmark}
\fancyhead[RO]{\small\rightmark}
\fancyfoot[C]{\small\thepage}
\renewcommand{\headrulewidth}{0.4pt}

% ----- Hiperlinks -----
\usepackage[colorlinks=true,linkcolor=NavyBlue,citecolor=ForestGreen,
            urlcolor=DarkRed,bookmarks=true]{hyperref}

% ----- Outros -----
\usepackage{enumitem}
\usepackage{float}
\usepackage{caption}
\usepackage{subcaption}
\usepackage{makeidx}
\makeindex

% ============================================================
%  AMBIENTES CUSTOMIZADOS
% ============================================================
\tcbset{
  defbox/.style={
    colback=blue!5,colframe=NavyBlue!70,
    fonttitle=\bfseries,breakable,
    before upper=\setlength{\parskip}{4pt}
  },
  teobox/.style={
    colback=green!4,colframe=ForestGreen!70,
    fonttitle=\bfseries,breakable
  },
  alertbox/.style={
    colback=red!5,colframe=DarkRed!70,
    fonttitle=\bfseries,breakable
  },
  notabox/.style={
    colback=yellow!8,colframe=Goldenrod!80,
    fonttitle=\bfseries,breakable
  }
}

\newtheorem{teorema}{Teorema}[chapter]
\newtheorem{proposicao}[teorema]{Proposição}
\newtheorem{lema}[teorema]{Lema}
\newtheorem{corolario}[teorema]{Corolário}
\theoremstyle{definition}
\newtheorem{definicao}{Definição}[chapter]
\newtheorem{exemplo}{Exemplo}[chapter]
\theoremstyle{remark}
\newtheorem{observacao}{Observação}[chapter]

% ============================================================
%  COMANDOS MATEMÁTICOS
% ============================================================
\newcommand{\R}{\mathbb{R}}
\newcommand{\E}{\mathbb{E}}
\newcommand{\Var}{\mathrm{Var}}
\newcommand{\Cov}{\mathrm{Cov}}
\newcommand{\Prob}{\mathbb{P}}
\newcommand{\N}{\mathcal{N}}
\newcommand{\D}{\mathcal{D}}
\newcommand{\Loss}{\mathcal{L}}
\newcommand{\bx}{\bm{x}}
\newcommand{\by}{\bm{y}}
\newcommand{\bw}{\bm{w}}
\newcommand{\bbeta}{\bm{\beta}}
\newcommand{\btheta}{\bm{\theta}}
\newcommand{\norm}[1]{\left\|#1\right\|}
\newcommand{\abs}[1]{\left|#1\right|}
\newcommand{\inner}[2]{\langle #1,\, #2 \rangle}
\DeclareMathOperator*{\argmin}{arg\,min}
\DeclareMathOperator*{\argmax}{arg\,max}
\DeclareMathOperator{\sign}{sign}
\DeclareMathOperator{\tr}{tr}
\DeclareMathOperator{\diag}{diag}
\DeclareMathOperator{\rank}{rank}
\DeclareMathOperator{\logit}{logit}

% ============================================================
%  INÍCIO DO DOCUMENTO
% ============================================================
\begin{document}

% ---- Capa ----
\begin{titlepage}
\centering
\vspace*{2cm}
{\Huge\bfseries Apostila de Machine Learning\par}
\vspace{0.5cm}
{\LARGE Teoria Avançada e Fundamentos Matemáticos\par}
\vspace{2cm}
\rule{\linewidth}{0.5pt}
\vspace{0.5cm}

{\large
\textbf{Módulos:}\\[6pt]
Introdução $\cdot$ Pré-processamento $\cdot$ Feature Engineering\\
Regressão $\cdot$ Classificação $\cdot$ Avaliação de Modelos\\
Regras de Associação $\cdot$ Clustering $\cdot$ Redes Neurais\\
SVM $\cdot$ Ensemble $\cdot$ Redes Bayesianas\\
Seleção de Variáveis $\cdot$ Interpretabilidade $\cdot$ Deploy $\cdot$ Armadilhas
}

\vspace{0.5cm}
\rule{\linewidth}{0.5pt}
\vfill
{\large Em português $\cdot$ Com exercícios e demonstrações\par}
\vspace{1cm}
{\large \today\par}
\end{titlepage}

\frontmatter
\tableofcontents
\listoffigures
\listoftables

% ============================================================
\mainmatter
% ============================================================

% ############################################################
\chapter{Introdução ao Machine Learning}
% ############################################################

\section{O que é Machine Learning?}

\begin{tcolorbox}[defbox,title=Definição Formal]
Seja $\D = \{(\bx_i, y_i)\}_{i=1}^n$ um conjunto de dados com entradas $\bx_i \in \R^p$ e saídas $y_i \in \mathcal{Y}$. O objetivo do aprendizado de máquina é encontrar uma função $f: \R^p \to \mathcal{Y}$ que minimize o \textbf{risco esperado}:
\[
R(f) = \E_{(\bx,y) \sim P}\bigl[\ell(f(\bx), y)\bigr]
\]
onde $\ell$ é uma função de perda e $P$ é a distribuição conjunta desconhecida dos dados.
\end{tcolorbox}

Como $P$ é desconhecida, minimizamos o \textbf{risco empírico}:
\[
\hat{R}(f) = \frac{1}{n}\sum_{i=1}^n \ell\bigl(f(\bx_i), y_i\bigr)
\]

\section{Tipos de Aprendizado}

\subsection{Aprendizado Supervisionado}

O modelo aprende um mapeamento $f: \mathcal{X} \to \mathcal{Y}$ a partir de pares $(\bx_i, y_i)$.

\begin{itemize}
  \item \textbf{Regressão}: $\mathcal{Y} = \R$ (saída contínua)
  \item \textbf{Classificação}: $\mathcal{Y} = \{1,\ldots,K\}$ (saída discreta)
\end{itemize}

\subsection{Aprendizado Não Supervisionado}

Apenas $\{\bx_i\}_{i=1}^n$ disponível, sem rótulos. O objetivo é descobrir a estrutura latente de $P(\bx)$.

\subsection{Aprendizado por Reforço}

Um agente interage com um ambiente em estados $s \in \mathcal{S}$, toma ações $a \in \mathcal{A}$ e recebe recompensas $r_t$. O objetivo é maximizar o retorno acumulado:
\[
G_t = \sum_{k=0}^{\infty} \gamma^k r_{t+k}, \quad \gamma \in [0,1)
\]

\section{Decomposição Viés-Variância}

Para qualquer estimador $\hat{f}$ treinado em $\D$:

\begin{tcolorbox}[teobox,title=Teorema da Decomposição Viés-Variância]
\[
\E_{\D}\bigl[(y - \hat{f}(\bx))^2\bigr] =
\underbrace{\bigl(\E_\D[\hat{f}(\bx)] - f(\bx)\bigr)^2}_{\text{Viés}^2}
+ \underbrace{\E_\D\bigl[(\hat{f}(\bx) - \E_\D[\hat{f}(\bx)])^2\bigr]}_{\text{Variância}}
+ \underbrace{\sigma^2_\epsilon}_{\text{Irredutível}}
\]
\end{tcolorbox}

\begin{itemize}
  \item \textbf{Viés alto} (underfitting): modelo muito simples.
  \item \textbf{Variância alta} (overfitting): modelo muito complexo, sensível ao conjunto de treino.
\end{itemize}

\section{Teoria de Aprendizado de Vapnik-Chervonenkis}

A \textbf{dimensão VC} ($d_{VC}$) de uma classe de hipóteses $\mathcal{H}$ mede sua capacidade expressiva. O limite de generalização de Vapnik-Chervonenkis afirma:

\[
R(f) \leq \hat{R}(f) + \sqrt{\frac{d_{VC}\left(\ln\frac{2n}{d_{VC}} + 1\right) + \ln\frac{4}{\delta}}{n}}
\]

com probabilidade $\geq 1 - \delta$. Isso formaliza o tradeoff entre complexidade do modelo e tamanho da amostra.

% ############################################################
\chapter{Pré-processamento de Dados}
% ############################################################

\section{Tratamento de Valores Ausentes}

Seja $X \in \R^{n \times p}$ a matriz de dados. Definimos o padrão de missingness pela matriz indicadora $M \in \{0,1\}^{n \times p}$, onde $M_{ij} = 1$ se $X_{ij}$ está ausente.

\paragraph{Mecanismos de ausência (Rubin, 1976):}
\begin{itemize}
  \item \textbf{MCAR} (Missing Completely at Random): $P(M \mid X_{\text{obs}}, X_{\text{mis}}) = P(M)$
  \item \textbf{MAR} (Missing at Random): $P(M \mid X_{\text{obs}}, X_{\text{mis}}) = P(M \mid X_{\text{obs}})$
  \item \textbf{MNAR} (Missing Not at Random): depende de $X_{\text{mis}}$
\end{itemize}

\subsection{Imputação por Máxima Verossimilhança (EM)}

O algoritmo EM trata os dados ausentes como variáveis latentes:

\textbf{E-step:} $Q(\theta \mid \theta^{(t)}) = \E_{X_{\text{mis}} \mid X_{\text{obs}}, \theta^{(t)}}[\log P(X \mid \theta)]$

\textbf{M-step:} $\theta^{(t+1)} = \argmax_\theta Q(\theta \mid \theta^{(t)})$

\section{Normalização e Padronização}

\subsection{Padronização Z-score}
\[
x'_j = \frac{x_j - \mu_j}{\sigma_j}, \quad \mu_j = \frac{1}{n}\sum_{i=1}^n x_{ij},\quad \sigma_j = \sqrt{\frac{1}{n-1}\sum_{i=1}^n(x_{ij}-\mu_j)^2}
\]

\subsection{Normalização Min-Max}
\[
x'_j = \frac{x_j - \min_j}{\max_j - \min_j} \in [0,1]
\]

\subsection{Normalização Robusta (Mediana/IQR)}

Resistente a outliers:
\[
x'_j = \frac{x_j - \tilde{x}_j}{Q_3 - Q_1}
\]
onde $\tilde{x}_j$ é a mediana e $Q_1, Q_3$ são os quartis.

\section{Detecção de Outliers}

\subsection{Regra de Tukey (IQR)}
\[
\text{outlier se } x < Q_1 - 1.5\,\mathrm{IQR} \quad \text{ou} \quad x > Q_3 + 1.5\,\mathrm{IQR}
\]

\subsection{Distância de Mahalanobis}

Para detectar outliers multivariados:
\[
D_M(\bx) = \sqrt{(\bx - \bm{\mu})^T \Sigma^{-1} (\bx - \bm{\mu})}
\]
Sob normalidade multivariada, $D_M^2 \sim \chi^2_p$. Valores com $p$-valor $< 0.001$ são candidatos a outliers.

\section{Weight of Evidence e Information Value}

Amplamente usados em scorecard de crédito:
\[
\mathrm{WOE}_i = \ln\!\left(\frac{p_{b,i}}{p_{m,i}}\right) = \ln\!\left(\frac{B_i/B_T}{M_i/M_T}\right)
\]
onde $B_i$ = bons no grupo $i$, $M_i$ = maus no grupo $i$, $B_T$ e $M_T$ são os totais.

\[
\mathrm{IV} = \sum_{i=1}^{k} (p_{b,i} - p_{m,i}) \cdot \mathrm{WOE}_i
\]

\begin{table}[H]
\centering
\caption{Interpretação do Information Value}
\begin{tabular}{cc}
\toprule
\textbf{IV} & \textbf{Poder Preditivo} \\
\midrule
$< 0.02$ & Inútil \\
$0.02$--$0.1$ & Fraco \\
$0.1$--$0.3$ & Médio \\
$0.3$--$0.5$ & Forte \\
$> 0.5$ & Suspeito (possível data leakage) \\
\bottomrule
\end{tabular}
\end{table}

% ############################################################
\chapter{Feature Engineering}
% ############################################################

\section{Transformações de Estabilização de Variância}

\subsection{Família Box-Cox}

Para $x > 0$:
\[
x^{(\lambda)} = \begin{cases}
\dfrac{x^\lambda - 1}{\lambda} & \text{se } \lambda \neq 0 \\[6pt]
\ln x & \text{se } \lambda = 0
\end{cases}
\]

O parâmetro $\lambda$ é estimado por máxima verossimilhança, assumindo normalidade após a transformação:
\[
\hat{\lambda} = \argmax_\lambda \left[ -\frac{n}{2}\ln\hat{\sigma}^2(\lambda) + (\lambda-1)\sum_{i=1}^n \ln x_i \right]
\]

\subsection{Família Yeo-Johnson}

Estende Box-Cox para $x \in \R$:
\[
x^{(\lambda)} = \begin{cases}
\dfrac{(x+1)^\lambda - 1}{\lambda} & x \geq 0,\; \lambda \neq 0 \\[4pt]
\ln(x+1) & x \geq 0,\; \lambda = 0 \\[4pt]
-\dfrac{(-x+1)^{2-\lambda}-1}{2-\lambda} & x < 0,\; \lambda \neq 2 \\[4pt]
-\ln(-x+1) & x < 0,\; \lambda = 2
\end{cases}
\]

\section{Análise de Componentes Principais (PCA)}

\begin{definicao}[PCA]
Dado $X \in \R^{n \times p}$ centrado, PCA encontra a projeção ortogonal de menor dimensão $q < p$ que maximiza a variância retida:
\[
\max_{W \in \R^{p \times q},\; W^TW = I_q} \tr(W^T \Sigma W)
\]
onde $\Sigma = \frac{1}{n-1}X^TX$ é a matriz de covariância amostral.
\end{definicao}

\begin{teorema}[Solução via Decomposição Espectral]
Os $q$ componentes principais são as $q$ autovetores de $\Sigma$ correspondentes aos $q$ maiores autovalores $\lambda_1 \geq \lambda_2 \geq \cdots \geq \lambda_q$.
A variância explicada pelo $k$-ésimo componente é:
\[
\text{VE}_k = \frac{\lambda_k}{\sum_{j=1}^p \lambda_j}
\]
\end{teorema}

\section{Embeddings Não-Lineares: t-SNE}

t-SNE (van der Maaten \& Hinton, 2008) minimiza a divergência KL entre distribuições de similaridade no espaço original e no espaço de baixa dimensão.

\textbf{Similaridade no espaço original} (Gaussiana):
\[
p_{j|i} = \frac{\exp(-\norm{\bx_i - \bx_j}^2 / 2\sigma_i^2)}{\sum_{k \neq i} \exp(-\norm{\bx_i - \bx_k}^2 / 2\sigma_i^2)}, \quad p_{ij} = \frac{p_{j|i} + p_{i|j}}{2n}
\]

\textbf{Similaridade no espaço 2D} (distribuição $t$ com 1 grau de liberdade — cauda pesada):
\[
q_{ij} = \frac{(1 + \norm{\by_i - \by_j}^2)^{-1}}{\sum_{k \neq l}(1 + \norm{\by_k - \by_l}^2)^{-1}}
\]

\textbf{Função objetivo:}
\[
\mathcal{C} = \mathrm{KL}(P \| Q) = \sum_{i \neq j} p_{ij} \ln \frac{p_{ij}}{q_{ij}}
\]

Minimizada por gradiente descendente. A \textbf{perplexidade} $\mathcal{P} = 2^{H(P_i)}$ controla o número efetivo de vizinhos.

\section{Features RFM}

Segmentação comportamental clássica em CRM:
\[
R = \text{dias desde última compra}, \quad F = \text{frequência de compras}, \quad M = \sum \text{valores gastos}
\]

O score RFM combinado: $\mathrm{RFM} = w_R \cdot \tilde{R} + w_F \cdot \tilde{F} + w_M \cdot \tilde{M}$, onde $\tilde{\cdot}$ indica percentil transformado em quintil $(1$--$5)$.

% ############################################################
\chapter{Regressão}
% ############################################################

\section{Mínimos Quadrados Ordinários}

\begin{definicao}[Modelo Linear]
$y = X\bbeta + \bm{\varepsilon}$, com $\bm{\varepsilon} \sim \N(\bm{0}, \sigma^2 I_n)$, $X \in \R^{n \times (p+1)}$ incluindo coluna de uns.
\end{definicao}

\begin{teorema}[Estimador OLS]
O estimador que minimiza $\mathrm{RSS} = \norm{\by - X\bbeta}^2$ é:
\[
\hat{\bbeta}_{\mathrm{OLS}} = (X^TX)^{-1}X^T\by
\]
\end{teorema}

\begin{teorema}[Gauss-Markov]
Sob as hipóteses clássicas de regressão, $\hat{\bbeta}_{\mathrm{OLS}}$ é o \textbf{melhor estimador linear não viesado} (BLUE): tem variância mínima na classe dos estimadores lineares não viesados.
\end{teorema}

A matriz de covariância do estimador OLS é:
\[
\Var(\hat{\bbeta}_{\mathrm{OLS}}) = \sigma^2 (X^TX)^{-1}, \quad \hat{\sigma}^2 = \frac{\mathrm{RSS}}{n-p-1}
\]

\section{Regularização}

\subsection{Ridge (Tikhonov)}

\[
\hat{\bbeta}_{\mathrm{Ridge}} = \argmin_\bbeta \norm{\by - X\bbeta}^2 + \lambda \norm{\bbeta}^2_2
\]

Solução fechada:
\[
\hat{\bbeta}_{\mathrm{Ridge}} = (X^TX + \lambda I)^{-1} X^T\by
\]

Via SVD ($X = U D V^T$):
\[
\hat{\bbeta}_{\mathrm{Ridge}} = \sum_{j=1}^p \frac{d_j^2}{d_j^2 + \lambda} \frac{\bm{u}_j^T\by}{d_j} \bm{v}_j
\]
Ridge contrai os coeficientes suavemente; multicolinearidade é controlada pois $(X^TX + \lambda I)$ é sempre invertível.

\subsection{LASSO}

\[
\hat{\bbeta}_{\mathrm{LASSO}} = \argmin_\bbeta \norm{\by - X\bbeta}^2 + \lambda \norm{\bbeta}_1
\]

Solução via \textit{soft-thresholding} para preditores ortogonais:
\[
\hat{\beta}_j^{\mathrm{LASSO}} = \sign(\hat{\beta}_j^{\mathrm{OLS}})\bigl(\abs{\hat{\beta}_j^{\mathrm{OLS}}} - \lambda/2\bigr)_+
\]
A penalidade $\ell_1$ gera soluções \textbf{esparsas} — coeficientes pequenos são zerados exatamente.

\subsection{Elastic Net}

\[
\hat{\bbeta}_{\mathrm{EN}} = \argmin_\bbeta \norm{\by - X\bbeta}^2 + \lambda_1 \norm{\bbeta}_1 + \lambda_2 \norm{\bbeta}^2_2
\]

Parametrização equivalente: $\alpha = \lambda_1/(\lambda_1 + \lambda_2) \in [0,1]$ (razão L1).

\section{Modelos Lineares Generalizados (GLM)}

\begin{definicao}[GLM]
Um GLM especifica:
\begin{enumerate}
  \item Distribuição da família exponencial: $p(y;\theta,\phi) = \exp\!\left[\frac{y\theta - b(\theta)}{a(\phi)} + c(y,\phi)\right]$
  \item Preditor linear: $\eta_i = \bx_i^T\bbeta$
  \item Função de ligação $g$: $g(\mu_i) = \eta_i$, onde $\mu_i = \E[y_i]$
\end{enumerate}
\end{definicao}

\begin{table}[H]
\centering
\caption{Componentes típicos de GLMs}
\begin{tabular}{llll}
\toprule
\textbf{Resposta} & \textbf{Distribuição} & \textbf{Ligação canônica} & \textbf{Aplicação} \\
\midrule
Contínua & Normal & Identidade $g(\mu)=\mu$ & Regressão linear \\
Binária & Binomial & Logit $g(\mu)=\ln\frac{\mu}{1-\mu}$ & Classificação \\
Contagem & Poisson & Log $g(\mu)=\ln\mu$ & Freq. de sinistros \\
Positiva & Gamma & Log ou inversa & Severidade de perdas \\
\bottomrule
\end{tabular}
\end{table}

Os parâmetros são estimados por máxima verossimilhança via algoritmo IRLS (\textit{Iteratively Reweighted Least Squares}):
\[
\bbeta^{(t+1)} = (X^T W^{(t)} X)^{-1} X^T W^{(t)} \bm{z}^{(t)}
\]
onde $W^{(t)} = \diag\!\left(\frac{(\partial\mu_i/\partial\eta_i)^2}{\Var(y_i)}\right)$ e $z_i^{(t)} = \eta_i^{(t)} + (y_i - \mu_i^{(t)})\frac{\partial\eta_i}{\partial\mu_i}$.

\section{Modelos Aditivos Generalizados (GAM)}

\[
g(\mu_i) = \alpha + f_1(x_{i1}) + f_2(x_{i2}) + \cdots + f_p(x_{ip})
\]

Cada $f_j$ é uma função suave estimada por \textit{splines} de suavização, encontrada via backfitting:
\[
\hat{f}_j \leftarrow \mathcal{S}_j\!\left[\by - \alpha - \sum_{k \neq j} \hat{f}_k(\bm{x}_k)\right]
\]
onde $\mathcal{S}_j$ é um operador de suavização. GAMs preservam a interpretabilidade (efeito parcial de cada variável) enquanto capturam não-linearidades.

% ############################################################
\chapter{Classificação}
% ############################################################

\section{Regressão Logística}

\subsection{Modelo e Estimação}

\[
P(y=1 \mid \bx) = \sigma(\bx^T\bbeta) = \frac{1}{1+e^{-\bx^T\bbeta}}
\]

Log-verossimilhança:
\[
\ell(\bbeta) = \sum_{i=1}^n \Bigl[y_i \log\sigma(\bx_i^T\bbeta) + (1-y_i)\log(1-\sigma(\bx_i^T\bbeta))\Bigr]
\]

Gradiente: $\nabla_\bbeta \ell = X^T(\by - \hat{\bm{p}})$

Hessiana: $H = -X^T W X$, onde $W = \diag(\hat{p}_i(1-\hat{p}_i))$

\subsection{Interpretação via Odds Ratio}

\[
\mathrm{OR}_j = e^{\beta_j}: \text{ aumento de 1 unidade em } x_j \text{ multiplica o odds por } e^{\beta_j}
\]

\subsection{Regularização em Classificação}

Equivalente a colocar priors nos coeficientes:
\begin{itemize}
  \item Prior Gaussiano $\Rightarrow$ penalidade Ridge (L2)
  \item Prior Laplaciano $\Rightarrow$ penalidade LASSO (L1)
\end{itemize}

\section{Árvores de Decisão}

\subsection{Critérios de Split}

\textbf{Índice de Gini}:
\[
G(t) = \sum_{k=1}^K p_k(t)(1-p_k(t)) = 1 - \sum_{k=1}^K p_k(t)^2
\]

\textbf{Entropia}:
\[
H(t) = -\sum_{k=1}^K p_k(t)\log_2 p_k(t)
\]

\textbf{Ganho de informação}:
\[
\mathrm{IG}(t, s) = H(t) - \frac{n_L}{n}H(t_L) - \frac{n_R}{n}H(t_R)
\]

\subsection{Poda (Pruning)}

\textbf{Minimal Cost-Complexity Pruning}: encontra a subárvore $T_\alpha$ que minimiza:
\[
R_\alpha(T) = R(T) + \alpha |T|
\]
onde $R(T)$ é o erro de resubstituição e $|T|$ é o número de folhas.

\section{K-Vizinhos Mais Próximos (KNN)}

\textbf{Classificação}: $\hat{y} = \mathrm{mode}\{y_j : j \in \mathcal{N}_K(\bx)\}$

\textbf{Métricas de distância}:
\begin{itemize}
  \item Euclidiana: $d_2(\bx,\bx') = \norm{\bx - \bx'}_2$
  \item Manhattan: $d_1(\bx,\bx') = \norm{\bx - \bx'}_1$
  \item Minkowski: $d_p(\bx,\bx') = \norm{\bx - \bx'}_p$
  \item Mahalanobis: $d_M(\bx,\bx') = \sqrt{(\bx-\bx')^T\Sigma^{-1}(\bx-\bx')}$
\end{itemize}

\textbf{Maldição da dimensionalidade}: em $p$ dimensões, a fração do volume total contida na hiperesfera de raio $r$ decresce como $(r/R)^p \to 0$.

% ############################################################
\chapter{Avaliação de Modelos}
% ############################################################

\section{Matriz de Confusão e Métricas Derivadas}

\begin{table}[H]
\centering
\caption{Matriz de confusão binária}
\begin{tabular}{lcc}
\toprule
 & \textbf{Predito Positivo} & \textbf{Predito Negativo} \\
\midrule
\textbf{Real Positivo} & VP (Verdadeiro Positivo) & FN (Falso Negativo) \\
\textbf{Real Negativo} & FP (Falso Positivo) & VN (Verdadeiro Negativo) \\
\bottomrule
\end{tabular}
\end{table}

\[
\text{Acurácia} = \frac{\mathrm{VP}+\mathrm{VN}}{n}, \quad
\text{Precisão} = \frac{\mathrm{VP}}{\mathrm{VP}+\mathrm{FP}}, \quad
\text{Recall} = \frac{\mathrm{VP}}{\mathrm{VP}+\mathrm{FN}}
\]

\[
F_\beta = (1+\beta^2)\frac{\text{Precisão}\cdot\text{Recall}}{\beta^2\cdot\text{Precisão}+\text{Recall}}, \quad
\text{MCC} = \frac{\mathrm{VP}\cdot\mathrm{VN} - \mathrm{FP}\cdot\mathrm{FN}}{\sqrt{(\mathrm{VP}+\mathrm{FP})(\mathrm{VP}+\mathrm{FN})(\mathrm{VN}+\mathrm{FP})(\mathrm{VN}+\mathrm{FN})}}
\]

\section{Curva ROC e AUC}

A curva ROC plota Sensibilidade vs.\ $(1-\text{Especificidade})$ para todos os thresholds $\tau$:
\[
\mathrm{TPR}(\tau) = P(\hat{p} > \tau \mid y=1), \quad \mathrm{FPR}(\tau) = P(\hat{p} > \tau \mid y=0)
\]

\begin{tcolorbox}[teobox,title=Interpretação probabilística da AUC]
\[
\mathrm{AUC} = P(\hat{p}(y=1) > \hat{p}(y=0)) = \frac{1}{n_+ n_-}\sum_{i: y_i=1}\sum_{j: y_j=0} \mathbbm{1}[\hat{p}_i > \hat{p}_j]
\]
É equivalente à estatística de Wilcoxon-Mann-Whitney.
\end{tcolorbox}

\section{Estatística KS e Curva CAP}

\textbf{KS} (Kolmogorov-Smirnov):
\[
\mathrm{KS} = \max_\tau \abs{F_{\text{pos}}(\tau) - F_{\text{neg}}(\tau)}
\]

\textbf{Accuracy Ratio} (Curva CAP):
\[
\mathrm{AR} = 2 \cdot \mathrm{AUC} - 1
\]

\section{Validação Cruzada}

\textbf{$k$-fold CV}: dividir $\D$ em $k$ blocos $\D_1,\ldots,\D_k$; para $j=1,\ldots,k$ treinar em $\D \setminus \D_j$ e avaliar em $\D_j$:
\[
\hat{R}_{\mathrm{CV}} = \frac{1}{k}\sum_{j=1}^k \hat{R}_j
\]

\textbf{Viés e variância da CV}: $k=n$ (leave-one-out) tem viés mínimo mas variância alta; $k=5$ ou $k=10$ equilibra os dois.

\textbf{Correção de Nadeau-Bengio} para CV repetido:
\[
\Var(\hat{R}_{\mathrm{CV}}) \approx \left(\frac{1}{k} + \frac{n_{\text{test}}}{n_{\text{train}}}\right)\sigma^2_\epsilon
\]

% ############################################################
\chapter{Regras de Associação}
% ############################################################

\section{Definições Formais}

Seja $\mathcal{I} = \{i_1,\ldots,i_m\}$ o conjunto de itens e $\mathcal{T} = \{t_1,\ldots,t_n\}$ a coleção de transações, com $t_j \subseteq \mathcal{I}$.

\begin{definicao}[Suporte, Confiança e Lift]
Para itemsets $A, B \subseteq \mathcal{I}$:
\[
\mathrm{supp}(A) = \frac{|\{t \in \mathcal{T}: A \subseteq t\}|}{|\mathcal{T}|}
\]
\[
\mathrm{conf}(A \Rightarrow B) = \frac{\mathrm{supp}(A \cup B)}{\mathrm{supp}(A)} = P(B \mid A)
\]
\[
\mathrm{lift}(A \Rightarrow B) = \frac{\mathrm{conf}(A \Rightarrow B)}{\mathrm{supp}(B)} = \frac{P(A \cup B)}{P(A)P(B)}
\]
\end{definicao}

\begin{tcolorbox}[defbox,title=Interpretação do Lift]
\begin{itemize}
  \item $\mathrm{lift} = 1$: $A$ e $B$ são independentes.
  \item $\mathrm{lift} > 1$: associação positiva — $A$ aumenta a probabilidade de $B$.
  \item $\mathrm{lift} < 1$: associação negativa.
\end{itemize}
\end{tcolorbox}

\textbf{Conviction}:
\[
\mathrm{conv}(A \Rightarrow B) = \frac{1 - \mathrm{supp}(B)}{1 - \mathrm{conf}(A \Rightarrow B)}
\]
Alta conviction indica que a regra seria muito imprecisa se $A$ e $B$ fossem independentes.

\section{Algoritmo Apriori}

\begin{tcolorbox}[defbox,title=Propriedade Anti-monotônica (Apriori)]
Se $A \not\supseteq \sigma_{\min}$, então $\forall B \supseteq A$: $\mathrm{supp}(B) \leq \mathrm{supp}(A) < \sigma_{\min}$.
\end{tcolorbox}

Esta propriedade permite podar o espaço de busca: nenhum superconjunto de um itemset infrequente precisa ser avaliado.

\textbf{Complexidade}: no pior caso $O(2^m)$ itemsets, mas a poda torna o algoritmo tratável na prática.

\section{Regras de Sequências}

Uma sequência $\langle s_1, s_2, \ldots, s_k \rangle$ é frequente se sua frequência (respeitando a ordem temporal) excede $\sigma_{\min}$. O algoritmo PrefixSpan cresce sequências por extensão de prefixo, evitando geração de candidatos.

% ############################################################
\chapter{Técnicas de Clustering}
% ############################################################

\section{K-means}

\begin{definicao}[Problema de Otimização]
\[
\min_{C_1,\ldots,C_K,\, \mu_1,\ldots,\mu_K} \sum_{k=1}^K \sum_{\bx \in C_k} \norm{\bx - \bm{\mu}_k}^2
\]
\end{definicao}

O problema é NP-difícil em geral; o algoritmo de Lloyd encontra um mínimo local:

\begin{algorithm}[H]
\caption{Algoritmo de Lloyd (K-means)}
\KwIn{$X$, $K$, centroides iniciais $\bm{\mu}_1^{(0)},\ldots,\bm{\mu}_K^{(0)}$}
\Repeat{sem mudança nos centroides}{
  \textbf{Atribuição}: $C_k^{(t)} \leftarrow \{\bx_i : k = \argmin_j \norm{\bx_i - \bm{\mu}_j^{(t-1)}}^2\}$\;
  \textbf{Atualização}: $\bm{\mu}_k^{(t)} \leftarrow \dfrac{1}{|C_k^{(t)}|}\displaystyle\sum_{\bx \in C_k^{(t)}} \bx$\;
}
\end{algorithm}

\textbf{Complexidade}: $O(nKpt)$, onde $t$ é o número de iterações.

\textbf{K-means++}: inicialização por amostragem proporcional à distância ao centroide mais próximo:
\[
P(\bx_i) = \frac{d(\bx_i, \mathcal{C})^2}{\sum_j d(\bx_j, \mathcal{C})^2}
\]
Garante solução em expectativa $O(\log K)$ vezes o ótimo.

\section{Critérios de Avaliação de Clustering}

\textbf{Coeficiente de Silhueta}:
\[
s(i) = \frac{b(i) - a(i)}{\max(a(i), b(i))}, \quad s(i) \in [-1, 1]
\]
onde $a(i)$ é a distância intra-cluster média e $b(i)$ é a menor distância inter-cluster média.

\textbf{Índice Davies-Bouldin}:
\[
\mathrm{DB} = \frac{1}{K}\sum_{k=1}^K \max_{j \neq k} \frac{s_k + s_j}{d(\bm{\mu}_k, \bm{\mu}_j)}
\]
Menor é melhor ($\mathrm{DB} = 0$ é perfeito).

\textbf{Índice de Calinski-Harabasz} (razão de variâncias entre/intra):
\[
\mathrm{CH} = \frac{\mathrm{tr}(B_K)/(K-1)}{\mathrm{tr}(W_K)/(n-K)}
\]
Maior é melhor.

\section{DBSCAN}

\begin{definicao}[$\varepsilon$-vizinhança e Conceitos Centrais]
\[
N_\varepsilon(\bx) = \{\bx' \in \D : d(\bx,\bx') \leq \varepsilon\}
\]
\begin{itemize}
  \item \textbf{Core point}: $|N_\varepsilon(\bx)| \geq \mathrm{MinPts}$
  \item \textbf{Border point}: não é core, mas está em $N_\varepsilon$ de um core point
  \item \textbf{Noise}: não é core nem border
\end{itemize}
\end{definicao}

\textbf{Vantagens}: descobre clusters de forma arbitrária; robusto a outliers; não requer $K$.

\textbf{Complexidade}: $O(n \log n)$ com índices espaciais.

\section{Clustering Hierárquico}

A distância entre clusters no método de Ward minimiza a variância total intra-cluster após fusão:
\[
d_{\mathrm{Ward}}(A, B) = \sqrt{\frac{2|A||B|}{|A|+|B|}} \norm{\bm{\mu}_A - \bm{\mu}_B}
\]

% ############################################################
\chapter{Redes Neurais Artificiais}
% ############################################################

\section{Perceptron e MLP}

\subsection{Neurônio Formal}

\[
z = \sum_{j=1}^p w_j x_j + b = \bw^T\bx + b, \qquad \hat{y} = \varphi(z)
\]

\subsection{Funções de Ativação}

\begin{align}
\text{Sigmoid: } &\sigma(z) = \frac{1}{1+e^{-z}}, \quad \sigma'(z) = \sigma(z)(1-\sigma(z)) \label{eq:sigmoid}\\[4pt]
\text{Tanh: } &\tau(z) = \frac{e^z - e^{-z}}{e^z + e^{-z}}, \quad \tau'(z) = 1 - \tau(z)^2 \\[4pt]
\text{ReLU: } &\rho(z) = \max(0,z), \quad \rho'(z) = \mathbbm{1}[z>0] \\[4pt]
\text{GELU: } &g(z) = z\,\Phi(z), \quad \Phi = \text{CDF da Normal}
\end{align}

\subsection{Backpropagation}

Seja $L = \ell(\hat{y}, y)$ a perda. Para uma rede com $\ell$ camadas, o gradiente se propaga pela regra da cadeia:

\[
\frac{\partial L}{\partial W^{(\ell)}} = \frac{\partial L}{\partial \bm{a}^{(\ell)}} \cdot \frac{\partial \bm{a}^{(\ell)}}{\partial \bm{z}^{(\ell)}} \cdot \frac{\partial \bm{z}^{(\ell)}}{\partial W^{(\ell)}}
\]

Definindo o \textbf{delta de erro} $\bm{\delta}^{(\ell)} = \nabla_{\bm{a}^{(\ell)}} L \odot \varphi'(\bm{z}^{(\ell)})$:

\[
\bm{\delta}^{(\ell)} = \bigl(W^{(\ell+1)}\bigr)^T \bm{\delta}^{(\ell+1)} \odot \varphi'(\bm{z}^{(\ell)})
\]

\[
\nabla_{W^{(\ell)}} L = \bm{\delta}^{(\ell)} \bigl(\bm{a}^{(\ell-1)}\bigr)^T, \quad \nabla_{b^{(\ell)}} L = \bm{\delta}^{(\ell)}
\]

\section{Otimização}

\textbf{SGD com momentum}:
\[
\bm{v}_{t+1} = \gamma \bm{v}_t + \eta \nabla_\theta L, \qquad \theta_{t+1} = \theta_t - \bm{v}_{t+1}
\]

\textbf{Adam} (Kingma \& Ba, 2015):
\[
m_t = \beta_1 m_{t-1} + (1-\beta_1) g_t, \quad v_t = \beta_2 v_{t-1} + (1-\beta_2) g_t^2
\]
\[
\hat{m}_t = \frac{m_t}{1-\beta_1^t}, \quad \hat{v}_t = \frac{v_t}{1-\beta_2^t}
\]
\[
\theta_{t+1} = \theta_t - \frac{\eta}{\sqrt{\hat{v}_t} + \varepsilon}\hat{m}_t
\]

\section{Regularização em Redes Neurais}

\textbf{Dropout}: durante o treino, cada neurônio é ``desligado'' com probabilidade $p$. Equivale aproximadamente a treinar um ensemble de $2^n$ redes.

\textbf{Batch Normalization}: para ativações de uma mini-batch $\{z_i\}$:
\[
\hat{z}_i = \frac{z_i - \mu_\mathcal{B}}{\sqrt{\sigma^2_\mathcal{B} + \varepsilon}}, \quad y_i = \gamma \hat{z}_i + \beta
\]
Reduz o problema de \textit{internal covariate shift}, permite taxas de aprendizado maiores.

\section{Teorema Universal de Aproximação}

\begin{teorema}[Cybenko, 1989; Hornik, 1991]
Uma MLP com uma camada oculta de largura suficiente e ativação contínua não-constante pode aproximar qualquer função contínua $f: [0,1]^p \to \R$ com precisão arbitrária $\varepsilon > 0$.
\end{teorema}

% ############################################################
\chapter{Support Vector Machines}
% ############################################################

\section{Classificador de Margem Máxima}

\begin{definicao}[Hiperplano separador]
$H = \{\bx : \bw^T\bx + b = 0\}$, com distância de $\bx_i$ ao hiperplano: $\frac{y_i(\bw^T\bx_i + b)}{\norm{\bw}}$.
\end{definicao}

\textbf{Problema primal (caso separável):}
\[
\min_{\bw, b} \frac{1}{2}\norm{\bw}^2 \quad \text{s.t.} \quad y_i(\bw^T\bx_i + b) \geq 1, \quad i=1,\ldots,n
\]

A margem geométrica é $2/\norm{\bw}$.

\textbf{Problema dual (via multiplicadores de Lagrange):}
\[
\max_{\bm{\alpha} \geq 0} \sum_i \alpha_i - \frac{1}{2}\sum_{i,j}\alpha_i\alpha_j y_i y_j \bx_i^T\bx_j \quad \text{s.t.} \quad \sum_i \alpha_i y_i = 0
\]

Os \textbf{vetores de suporte} são os pontos com $\alpha_i > 0$.

\section{Soft Margin (Caso Não Separável)}

\[
\min_{\bw,b,\bm{\xi}} \frac{1}{2}\norm{\bw}^2 + C\sum_{i=1}^n \xi_i \quad \text{s.t.} \quad y_i(\bw^T\bx_i+b) \geq 1-\xi_i,\; \xi_i \geq 0
\]

$C$ controla o tradeoff entre margem máxima e violações.

\section{Kernels e o Truque do Kernel}

\begin{definicao}[Kernel de Mercer]
Uma função $k: \mathcal{X} \times \mathcal{X} \to \R$ é um kernel válido se e somente se para qualquer $\{x_1,\ldots,x_n\}$, a matriz de Gram $K_{ij} = k(\bx_i, \bx_j)$ é simétrica positiva semi-definida.
\end{definicao}

Substituindo $\bx_i^T\bx_j$ por $k(\bx_i,\bx_j)$ no problema dual, opera-se em espaços de features de dimensão possivelmente infinita sem computá-los explicitamente.

\begin{table}[H]
\centering
\caption{Kernels comuns}
\begin{tabular}{lll}
\toprule
\textbf{Kernel} & \textbf{Fórmula} & \textbf{Parâmetros} \\
\midrule
Linear & $\bx^T\bx'$ & --- \\
Polinomial & $(\gamma\,\bx^T\bx' + r)^d$ & $\gamma, r, d$ \\
RBF/Gaussiano & $\exp(-\gamma\norm{\bx-\bx'}^2)$ & $\gamma$ \\
Sigmoide & $\tanh(\gamma\,\bx^T\bx' + r)$ & $\gamma, r$ \\
\bottomrule
\end{tabular}
\end{table}

\section{SVM para Regressão (SVR)}

\[
\min_{\bw,b,\xi,\xi^*} \frac{1}{2}\norm{\bw}^2 + C\sum_i(\xi_i + \xi_i^*)
\]
\[
\text{s.t.}\quad y_i - \bw^T\bx_i - b \leq \varepsilon + \xi_i, \quad \bw^T\bx_i + b - y_i \leq \varepsilon + \xi_i^*, \quad \xi_i, \xi_i^* \geq 0
\]

O tubo $\varepsilon$-insensível define uma zona de tolerância ao erro.

% ############################################################
\chapter{Métodos Ensemble}
% ############################################################

\section{Teoria dos Ensembles}

\begin{teorema}[Redução de Erro por Votação]
Se $m$ classificadores têm erro $\varepsilon < 0.5$ e são independentes, o erro do ensemble por voto majoritário é:
\[
\varepsilon_{\text{ensemble}} = \sum_{k=\lceil m/2 \rceil}^{m} \binom{m}{k} \varepsilon^k (1-\varepsilon)^{m-k} \ll \varepsilon
\]
\end{teorema}

\section{Bagging}

\textbf{Bootstrap Aggregating} (Breiman, 1996): treina $T$ modelos em amostras bootstrap $\D_1^*,\ldots,\D_T^*$.

\[
\hat{f}_{\text{bag}}(\bx) = \frac{1}{T}\sum_{t=1}^T \hat{f}_t(\bx) \quad \text{(regressão)}
\]

\textbf{Out-of-Bag Error}: cada observação está ausente de $\approx 36.8\%$ das amostras bootstrap ($\lim_{n\to\infty}(1-1/n)^n = e^{-1}$), provendo uma estimativa de generalização gratuita.

\section{Random Forest}

Adiciona aleatoriedade ao Bagging: em cada split, apenas um subconjunto de $m \approx \sqrt{p}$ features (classificação) ou $p/3$ (regressão) é considerado. Isso descorrelaciona as árvores, reduzindo a variância do ensemble.

\textbf{Importância por MDI} (\textit{Mean Decrease Impurity}):
\[
\mathrm{Imp}(X_j) = \frac{1}{T}\sum_{t=1}^T \sum_{v \in \mathcal{V}_j^t} \frac{n_v}{n}\,\Delta i(v)
\]
onde $\mathcal{V}_j^t$ são os nós da árvore $t$ que usam $X_j$, e $\Delta i(v)$ é a redução de impureza no nó $v$.

\section{Boosting}

\subsection{AdaBoost}

Cada iteração $t$ ajusta os pesos para focar nos exemplos errados:

\[
\varepsilon_t = \sum_{i=1}^n w_i^{(t)} \mathbbm{1}[y_i \neq h_t(\bx_i)]
\]
\[
\alpha_t = \frac{1}{2}\ln\!\left(\frac{1-\varepsilon_t}{\varepsilon_t}\right)
\]
\[
w_i^{(t+1)} = \frac{w_i^{(t)} \exp(-\alpha_t y_i h_t(\bx_i))}{Z_t}
\]

Predição final: $H(\bx) = \sign\!\left(\sum_t \alpha_t h_t(\bx)\right)$.

\begin{teorema}[Margem e Generalização — Schapire et al., 1998]
O erro de treino de AdaBoost decresce exponencialmente: após $T$ rounds,
\[
\text{Erro treino} \leq \exp\!\left(-2\sum_{t=1}^T \gamma_t^2\right), \quad \gamma_t = \frac{1}{2} - \varepsilon_t
\]
\end{teorema}

\subsection{Gradient Boosting}

\begin{definicao}[Gradient Boosting]
Em cada iteração $t$, ajusta-se um modelo fraco $h_t$ nos pseudo-resíduos:
\[
r_i^{(t)} = -\left.\frac{\partial \ell(y_i, \hat{y}_i)}{\partial \hat{y}_i}\right|_{\hat{y}_i = F_{t-1}(\bx_i)}
\]
\[
F_t(\bx) = F_{t-1}(\bx) + \eta \cdot h_t(\bx)
\]
\end{definicao}

\subsection{XGBoost}

Objetivo regularizado na iteração $t$:
\[
\mathcal{L}^{(t)} = \sum_{i=1}^n \ell(y_i, \hat{y}_i^{(t-1)} + f_t(\bx_i)) + \Omega(f_t)
\]
\[
\Omega(f_t) = \gamma T + \frac{\lambda}{2}\sum_{j=1}^T w_j^2
\]

Aproximação de Taylor de segunda ordem:
\[
\mathcal{L}^{(t)} \approx \sum_{i=1}^n \left[g_i f_t(\bx_i) + \frac{1}{2}h_i f_t(\bx_i)^2\right] + \Omega(f_t)
\]

Para uma folha $j$, o score ótimo é $w_j^* = -\frac{\sum_{i \in I_j} g_i}{\sum_{i \in I_j} h_i + \lambda}$.

O ganho de split:
\[
\text{Gain} = \frac{1}{2}\left[\frac{G_L^2}{H_L+\lambda} + \frac{G_R^2}{H_R+\lambda} - \frac{(G_L+G_R)^2}{H_L+H_R+\lambda}\right] - \gamma
\]

% ############################################################
\chapter{Redes Bayesianas e Classificadores Probabilísticos}
% ############################################################

\section{Fundamentos Bayesianos}

\begin{tcolorbox}[teobox,title=Teorema de Bayes]
\[
P(\theta \mid \D) = \frac{P(\D \mid \theta)\, P(\theta)}{P(\D)} = \frac{\text{Verossimilhança} \times \text{Prior}}{\text{Evidência}}
\]
\end{tcolorbox}

\textbf{Atualização sequencial}: com observações $x_1, x_2, \ldots$ chegando em sequência, o posterior da iteração $t$ serve como prior da iteração $t+1$:
\[
P(\theta \mid x_1,\ldots,x_t) \propto P(x_t \mid \theta)\, P(\theta \mid x_1,\ldots,x_{t-1})
\]

\section{Naive Bayes}

Sob a hipótese de independência condicional das features dada a classe:
\[
P(C_k \mid \bx) \propto P(C_k)\prod_{j=1}^p P(x_j \mid C_k)
\]

\textbf{Gaussian NB}: $P(x_j \mid C_k) = \N(x_j;\, \mu_{jk}, \sigma_{jk}^2)$

\textbf{Multinomial NB} (para texto): $P(\text{doc} \mid C_k) = \frac{n!}{\prod_j n_j!}\prod_j \theta_{jk}^{n_j}$

\textbf{Laplace smoothing} para evitar probabilidade zero:
\[
\hat{\theta}_{jk} = \frac{n_{jk} + \alpha}{n_k + \alpha p}
\]

\section{Estrutura de Redes Bayesianas}

\begin{definicao}[Rede Bayesiana]
Um DAG $G = (V, E)$ sobre variáveis aleatórias $X_1,\ldots,X_p$ define a fatorização:
\[
P(X_1,\ldots,X_p) = \prod_{i=1}^p P(X_i \mid \mathrm{Pa}(X_i))
\]
\end{definicao}

\textbf{D-separação}: $X \perp\!\!\!\perp Y \mid Z$ se toda trilha entre $X$ e $Y$ é bloqueada por $Z$.

% ############################################################
\chapter{Seleção de Variáveis}
% ############################################################

\section{Métodos de Filtro}

\subsection{Correlação de Spearman e Cramer's V}

Para variáveis categóricas:
\[
V = \sqrt{\frac{\chi^2/n}{\min(r-1,c-1)}}
\]

\subsection{Teste F de Fisher (ANOVA)}

\[
F(X_j) = \frac{\sum_k n_k(\bar{x}_{jk} - \bar{x}_j)^2 / (K-1)}{\sum_k \sum_{i \in C_k}(x_{ij} - \bar{x}_{jk})^2 / (n-K)}
\]

\subsection{Informação Mútua}

\[
I(X_j; Y) = \sum_{x,y} p(x,y)\log\frac{p(x,y)}{p(x)p(y)} = H(Y) - H(Y \mid X_j)
\]

\section{Métodos Wrapper}

\textbf{RFE} (Recursive Feature Elimination): treina o modelo, remove a feature de menor importância, repete.

\textbf{Forward Selection por AIC/BIC}:
\[
\mathrm{AIC} = 2k - 2\ln\hat{L}, \qquad \mathrm{BIC} = k\ln n - 2\ln\hat{L}
\]

BIC penaliza mais fortemente a complexidade e é consistente (seleciona o modelo verdadeiro com $n \to \infty$).

\section{Métodos Embutidos}

\textbf{LASSO} realiza seleção e estimação simultaneamente. O caminho de regularização é computado eficientemente pelo algoritmo LARS.

\textbf{Importância por permutação}:
\[
\mathrm{PI}(X_j) = \hat{R}(f, X_{\mathrm{perm}_j}) - \hat{R}(f, X)
\]

% ############################################################
\chapter{Interpretabilidade de Modelos (XAI)}
% ############################################################

\section{Explicações Globais}

\subsection{Partial Dependence Plot (PDP)}

\[
\hat{f}_S(\bx_S) = \E_{X_C}\bigl[\hat{f}(\bx_S, X_C)\bigr] \approx \frac{1}{n}\sum_{i=1}^n \hat{f}(\bx_S, \bx_{C,i})
\]

\subsection{Individual Conditional Expectation (ICE)}

Análogo ao PDP mas por observação, revelando heterogeneidade:
\[
\hat{f}^{(i)}_S(x_s) = \hat{f}(x_s, \bx_{C,i}), \quad i=1,\ldots,n
\]

O PDP é a média das curvas ICE.

\section{Explicações Locais}

\subsection{LIME}

\begin{definicao}[LIME — Ribeiro et al., 2016]
Para uma predição em $\bx$, LIME encontra uma explicação $g \in \mathcal{G}$ (modelo interpretável) que:
\[
\xi(\bx) = \argmin_{g \in \mathcal{G}} \mathcal{L}(f, g, \pi_\bx) + \Omega(g)
\]
onde $\pi_\bx(\bx') = \exp(-d(\bx,\bx')^2/\sigma^2)$ é um kernel de proximidade.
\end{definicao}

\subsection{Valores de Shapley}

\begin{definicao}[Shapley Value]
Para um jogador $j$ em um jogo de coalizão $v: 2^{[p]} \to \R$:
\[
\phi_j(v) = \sum_{S \subseteq [p] \setminus \{j\}} \frac{|S|!\,(p-|S|-1)!}{p!}\bigl[v(S \cup \{j\}) - v(S)\bigr]
\]
\end{definicao}

\begin{tcolorbox}[teobox,title=Axiomas dos Valores de Shapley]
São a \textbf{única} alocação satisfazendo:
\begin{enumerate}
  \item \textbf{Eficiência}: $\sum_j \phi_j = v([p]) - v(\emptyset)$ (soma à predição total)
  \item \textbf{Simetria}: features equivalentes têm o mesmo $\phi$
  \item \textbf{Nulidade}: features sem efeito têm $\phi_j = 0$
  \item \textbf{Aditividade}: $\phi_j(v+w) = \phi_j(v) + \phi_j(w)$
\end{enumerate}
\end{tcolorbox}

\textbf{SHAP TreeExplainer}: calcula Shapley values exatamente em $O(TLD^2)$ para tree ensembles (Chen et al., 2020), onde $T$ = árvores, $L$ = folhas, $D$ = profundidade.

% ############################################################
\chapter{Deploy e Governança de Modelos}
% ############################################################

\section{Population Stability Index (PSI)}

Detecta desvio na distribuição de entrada entre treino e produção:
\[
\mathrm{PSI} = \sum_{i=1}^{k} (\%A_i - \%E_i) \cdot \ln\!\left(\frac{\%A_i}{\%E_i}\right)
\]

\begin{table}[H]
\centering
\caption{Interpretação do PSI}
\begin{tabular}{cl}
\toprule
\textbf{PSI} & \textbf{Diagnóstico} \\
\midrule
$< 0.1$ & Estável — sem ação necessária \\
$0.1$--$0.2$ & Alerta — investigar causas \\
$> 0.2$ & Instável — retreinar o modelo \\
\bottomrule
\end{tabular}
\end{table}

\section{Fairness e Viés Algorítmico}

\textbf{Paridade Demográfica} (Statistical Parity):
\[
P(\hat{y}=1 \mid A=a) = P(\hat{y}=1 \mid A=b)
\]

\textbf{Equalized Odds} (Hardt et al., 2016):
\[
P(\hat{y}=1 \mid A=a, y=c) = P(\hat{y}=1 \mid A=b, y=c), \quad c \in \{0,1\}
\]

\textbf{Calibração}: $P(y=1 \mid \hat{p} = s) = s$ para todo $s \in [0,1]$.

\begin{tcolorbox}[alertbox,title=Impossibilidade de Fairness]
Chouldechova (2017) e Kleinberg et al. (2017) demonstraram que paridade demográfica, equalized odds e calibração são \textbf{simultaneamente impossíveis} quando as taxas de base diferem entre grupos. A escolha de critério depende do contexto ético e legal.
\end{tcolorbox}

\section{Monitoramento de Modelos em Produção}

Um framework de monitoramento robusto acompanha:

\begin{enumerate}
  \item \textbf{Data drift}: mudança em $P(\bx)$ — detectada via PSI, KS, MMD
  \item \textbf{Concept drift}: mudança em $P(y \mid \bx)$ — detectada via CUSUM, ADWIN
  \item \textbf{Performance metrics}: degradação de AUC, F1 em dados recentes com rótulos
  \item \textbf{Operational metrics}: latência, throughput, erros de infraestrutura
\end{enumerate}

\textbf{Maximum Mean Discrepancy (MMD)} para detectar mudança de distribuição:
\[
\mathrm{MMD}^2(P,Q) = \E_{x,x'\sim P}[k(x,x')] - 2\E_{x\sim P,y\sim Q}[k(x,y)] + \E_{y,y'\sim Q}[k(y,y')]
\]

% ############################################################
\chapter{Armadilhas e Boas Práticas em ML}
% ############################################################

\section{Data Leakage}

Ocorre quando informação do futuro (ou do conjunto de teste) contamina o treinamento. Tipos:

\begin{itemize}
  \item \textbf{Target leakage}: features calculadas após o evento que o modelo tenta predizer
  \item \textbf{Train-test contamination}: pré-processamento (normalização, imputação) ajustado em todo o dataset antes do split
\end{itemize}

\begin{tcolorbox}[alertbox,title=Regra de Ouro: Sempre use Pipelines]
Qualquer transformação que ``aprende'' parâmetros (média, desvio-padrão, componentes PCA, bins para WOE) deve ser ajustada \textbf{somente} nos dados de treino e apenas \textbf{aplicada} nos dados de teste/validação.
\end{tcolorbox}

\section{Sobreajuste e Capacidade do Modelo}

Para um polinômio de grau $d$ com $n$ pontos, o modelo tem $d+1$ parâmetros. O risco de sobreajuste cresce quando $d/n \to 1$.

A \textbf{complexidade efetiva} é melhor medida pelo número efetivo de parâmetros (em modelos regularizados):
\[
\mathrm{df}(\lambda) = \tr\!\bigl[X(X^TX + \lambda I)^{-1}X^T\bigr] = \sum_{j=1}^p \frac{d_j^2}{d_j^2 + \lambda}
\]

\section{Desbalanceamento de Classes}

Para classe positiva com proporção $\pi \ll 0.5$:

\textbf{SMOTE} (Chawla et al., 2002): gera exemplos sintéticos por interpolação entre vizinhos da classe minoritária:
\[
\bx_{\text{new}} = \bx_i + \delta \cdot (\bx_{nn} - \bx_i), \quad \delta \sim U(0,1)
\]

\textbf{Calibração de threshold}: ao invés de usar $\tau = 0.5$, otimize $\tau$ pelo F1, ou use:
\[
\tau^* = \frac{\pi_0 \cdot C_{FN}}{\pi_0 \cdot C_{FN} + (1-\pi_0) \cdot C_{FP}}
\]
onde $C_{FN}$ e $C_{FP}$ são os custos de erros tipo I e II.

\section{Múltiplos Testes e p-Hacking}

Ao realizar $m$ testes com nível $\alpha$, a taxa de erro família (FWER) sem correção cresce:
\[
\mathrm{FWER} = 1 - (1-\alpha)^m
\]

\textbf{Correção de Bonferroni}: $\alpha_{\mathrm{corr}} = \alpha/m$ (controla FWER, conservadora)

\textbf{Benjamini-Hochberg}: controla a \textit{False Discovery Rate} $\mathrm{FDR} = \E[\mathrm{FP}/\max(\mathrm{R},1)]$, menos conservadora.

\section{Maldição da Dimensionalidade}

Para pontos uniformemente distribuídos em $[0,1]^p$, a distância entre o ponto mais próximo e mais distante converge:
\[
\lim_{p\to\infty} \frac{\max_i d(\bx_0, \bx_i) - \min_i d(\bx_0, \bx_i)}{\min_i d(\bx_0, \bx_i)} = 0
\]

Todos os pontos ficam à mesma distância — o conceito de ``vizinho próximo'' perde significado.

O volume de uma hiperesfera de raio $r$ em dimensão $p$:
\[
V_p(r) = \frac{\pi^{p/2}}{\Gamma(p/2+1)} r^p
\]
A razão $V_p(r)/V_p(R) = (r/R)^p \to 0$ exponencialmente.

% ============================================================
%  APÊNDICES
% ============================================================
\appendix

\chapter{Fundamentos de Álgebra Linear}

\section{Decomposições Matriciais}

\textbf{Decomposição em Valores Singulares (SVD)}:
\[
X = U D V^T, \quad U \in \R^{n \times r},\; D = \diag(d_1,\ldots,d_r),\; V \in \R^{p \times r}
\]

com $d_1 \geq d_2 \geq \cdots \geq d_r > 0$ (valores singulares).

\textbf{Decomposição Espectral}: para $A$ simétrica:
\[
A = Q \Lambda Q^T, \quad Q = [\bm{q}_1 \cdots \bm{q}_p],\; \Lambda = \diag(\lambda_1,\ldots,\lambda_p)
\]

\section{Traço, Determinante e Normas}

\[
\tr(AB) = \tr(BA), \quad \det(AB) = \det(A)\det(B)
\]
\[
\norm{A}_F = \sqrt{\tr(A^TA)} = \sqrt{\sum_{i,j} a_{ij}^2}, \quad \norm{A}_2 = \sigma_{\max}(A)
\]

\chapter{Fundamentos de Probabilidade e Estatística}

\section{Distribuições Importantes}

\textbf{Família Exponencial}:
\[
p(y;\eta) = h(y)\exp\bigl(\eta^T T(y) - A(\eta)\bigr)
\]
Propriedades: $\E[T(y)] = \nabla A(\eta)$, $\Var[T(y)] = \nabla^2 A(\eta)$.

\section{Divergências}

\textbf{KL divergence} (não-simétrica):
\[
\mathrm{KL}(P\|Q) = \sum_x P(x)\ln\frac{P(x)}{Q(x)} \geq 0
\]

\textbf{Jensen-Shannon} (simétrica, limitada em $[0,\ln 2]$):
\[
\mathrm{JS}(P\|Q) = \frac{1}{2}\mathrm{KL}(P\|M) + \frac{1}{2}\mathrm{KL}(Q\|M), \quad M = \frac{P+Q}{2}
\]

\chapter{Referências e Leituras Complementares}

\begin{itemize}[leftmargin=1.5cm]
  \item Hastie, T., Tibshirani, R., Friedman, J. \textit{The Elements of Statistical Learning}. 2ª ed., Springer, 2009.
  \item Bishop, C.M. \textit{Pattern Recognition and Machine Learning}. Springer, 2006.
  \item Murphy, K.P. \textit{Machine Learning: A Probabilistic Perspective}. MIT Press, 2012.
  \item Goodfellow, I., Bengio, Y., Courville, A. \textit{Deep Learning}. MIT Press, 2016.
  \item Molnar, C. \textit{Interpretable Machine Learning}. 2022. \url{https://christophm.github.io/interpretable-ml-book/}
  \item Breiman, L. \textit{Random Forests}. Machine Learning, 45(1):5--32, 2001.
  \item Chen, T., Guestrin, C. \textit{XGBoost: A Scalable Tree Boosting System}. KDD, 2016.
  \item Lundberg, S., Lee, S.I. \textit{A Unified Approach to Interpreting Model Predictions}. NeurIPS, 2017.
  \item Ribeiro, M.T., Singh, S., Guestrin, C. \textit{Why Should I Trust You? Explaining the Predictions of Any Classifier}. KDD, 2016.
  \item Chawla, N.V. et al. \textit{SMOTE: Synthetic Minority Over-sampling Technique}. JAIR, 16:321--357, 2002.
\end{itemize}

% ============================================================
\backmatter
\printindex
% ============================================================

\end{document}
